\subsection{Scope and Limitations of the Methodology}

The methodology developed in this chapter has a limited and specific scope. It defines a physically motivated framework for directional thermoelastic wave prediction, but it does not claim full validation of real geological wave propagation. Real geological media are heterogeneous, fractured, porous, anisotropic, pressure-dependent, temperature-dependent, and often influenced by fluids. A complete description of such media would require more detailed geological characterization, more advanced constitutive laws, and extensive validation against field or laboratory data.

The system should therefore be interpreted as a research prototype rather than as an industrial-scale geophysical simulator. The model outputs are comparative directional predictions, not field-validated measurements of thermoelastic wave propagation. They are useful for examining trends, comparing model behavior, and connecting physical variables with predicted response quantities, but they should not be used as definitive predictions for real geological operations without further validation.

Another limitation is the simplified material representation. In the mathematical formulation, each rock type is represented by effective locally homogeneous and isotropic parameters. This is necessary for a clear bachelor-level formulation and for a consistent input-output framework. However, real rocks may contain bedding, fractures, pore networks, mineral fabrics, and stress-dependent cracks that cannot be fully represented by a single set of scalar effective parameters.

The directional representation is also simplified. Source--probe direction, distance, and azimuth provide a useful geometric description, but they do not fully capture all complexities of wave propagation in heterogeneous media. Reflection, refraction, scattering, attenuation, anisotropic phase velocity, fracture interaction, and local thermal stress concentrations may require more detailed modelling than the present methodology provides.

The model comparison framework is limited by the available implementation and by the available data or generated response structure. PINN, FNO, MeshGraphNet, and Transformer-style components are compared as alternative predictive approaches, but the comparison should not be interpreted as a universal ranking of these model families for all thermoelastic or geophysical problems. Their behavior depends on data, architecture, training assumptions, physical constraints, and implementation details.

Despite these limitations, the methodology provides a coherent bridge between the physical theory and the implemented prototype. Chapter 1 explained the physical and geological basis of thermoelastic wave propagation. Chapter 2 expressed that basis in mathematical variables, equations, parameters, and directional quantities. Chapter 3 uses the same structure to define a prediction methodology in which the physical fields and material parameters become inputs and outputs of a directional thermoelastic prediction framework. This prepares the thesis for the next chapter, where the COMSOL simulation setup is described.

\subsection{Summary of Chapter 3}

This chapter defined the methodology used to convert the thermoelastic formulation into a directional prediction task. The chapter connected the physical fields, material parameters, source position, observation point, distance, and azimuth introduced earlier with a common input--output structure for model comparison.

The chapter also clarified the interpretation of the proposed framework. The methodology is a research-prototype approach for comparing predictive components under the same physical scenario, not a claim of complete field validation. Its main value is that it keeps the comparison logically tied to the thermoelastic problem while making the limitations of simplified material representation, directional geometry, and available model implementations explicit.
